% ECONOMICS_PROBLEM: {"id": "O25-370", "title": "Industrial policy under weak institutions", "jel_code": "O25", "source_id": "OP-0033"}
% !TeX program = xelatex
\documentclass[11pt,a4paper]{article}
\usepackage[margin=23mm]{geometry}
\usepackage{fontspec}
\setmainfont{Latin Modern Roman}
\usepackage{amsmath,amssymb,mathtools}
\usepackage{xcolor,enumitem,booktabs,longtable,array,xurl,hyperref}
\definecolor{muted}{HTML}{53636C}
\hypersetup{colorlinks=true,urlcolor=blue,linkcolor=blue}
\setlength{\parindent}{0pt}
\setlength{\parskip}{5pt}
\newcommand{\R}{\mathbb R}
\newcommand{\E}{\mathbb E}
\newcommand{\N}{\mathbb N}
\newcommand{\ind}{\mathbf 1}
\newcommand{\Var}{\operatorname{Var}}
\newcommand{\Cov}{\operatorname{Cov}}
\newcommand{\supp}{\operatorname{supp}}
\DeclareMathOperator*{\argmin}{arg\,min}
\DeclareMathOperator*{\argmax}{arg\,max}
\newcommand{\entrymeta}[3]{{\sffamily\small\color{muted}\textbf{\texttt{#1}}\quad|\quad #2\par #3\par}}
\newcommand{\Problemref}[1]{\texttt{#1}}
\newcommand{\JELref}[2]{\texttt{#2} (JEL \texttt{#1})}
\begin{document}
\section*{Industrial policy under weak institutions}
\textbf{Problem ID:} \texttt{O25-370}\quad \textbf{JEL:} \texttt{O25}\par
% BEGIN ECONOMICS STATEMENT
\label{problem:OP-0033}
\label{atlas:prob:D3}
\noindent\textbf{Original atlas ID:} D3 (entry 33). \textbf{Classification:} Editorial robust industrial-policy design problem.

\noindent\textbf{Original atlas question:} When can targeted subsidies/credit accelerate structural transformation rather than invite capture?

\subsection*{Definitions and assumptions}
Fix sectors, firms, and horizon $H$. Let a government policy $\pi$ specify subsidies, credit guarantees, procurement, eligibility, information disclosure, and sunset rules as functions of observed histories. In model $M$, firms have private productivity and investment information; government signals are explicitly defined. Production spillovers, learning, input-output links, financing frictions, lobbying, and political replacement are included only when their laws are stated. Let $\mathcal P_M$ be the set of budget-feasible and incentive-compatible implementable policies. Specify a dynamic equilibrium and a selection rule. Let $W_M(\pi)$ be expected discounted household-and-owner welfare net of real administrative and financing costs, with announced distributional weights.
\subsection*{Formal research question}
Given observations $P$ and maintained restrictions, let $\mathcal M(P)$ be the nonempty set of compatible models. Determine whether a targeted policy can be shown to improve welfare relative to a specified untargeted baseline $\pi_0$ feasible in every compatible model:
\[
 \exists\pi\in\bigcap_{M\in\mathcal M(P)}\mathcal P_M
 \quad\text{such that}\quad
 \inf_{M\in\mathcal M(P)}[W_M(\pi)-W_M(\pi_0)]>0.
\]
When this robust condition fails, characterize model-specific welfare gains and their sensitivity to spillovers, capture, and government information. If the relevant supremum is finite, construct approximately optimal feasible rules and identify the observations required to distinguish their expected outcomes. Success is thus a conditional welfare statement, not a universal sign claim about subsidies.
\subsection*{Interpretation and scope}
A subsidy can raise a recipient's output while reducing aggregate productivity through displacement, financing, or inefficient entry. Treatment effects on recipients therefore differ from policy effects on the entire production network. Transfers are not real resource costs by themselves, but distortionary financing and administration can be costly. Learning spillovers must be distinguished from market-price spillovers and privately appropriated returns. Weak institutions should be represented by measurable information and implementation restrictions or a political model, rather than an unspecified adjective. Ex ante trials require credible treatment assignment and measurement of competitors, suppliers, consumers, and future political compliance. A nominal sunset does not imply an enforceable commitment to terminate support.
\subsection*{Source audit and status}
[S1] is Rodrik's 2004 working paper; [S2] a development-handbook review; [S3] the 2024 review of new industrial-policy economics. They verify the literature and motivate conditional design questions, while not stating the exact robust optimization criterion above. The formulation is editorial. Particular policy successes or failures remain established evidence, and universal 2026 unresolved status is not inferred from a review.

\subsection*{References}
\begin{enumerate}[label=\textnormal{[S\arabic*]},leftmargin=*]
\item Dani Rodrik (2004). \emph{Industrial Policy for the Twenty-First Century}. Harvard Kennedy School Faculty Research Working Paper RWP04-047. \url{https://www.hks.harvard.edu/publications/industrial-policy-twenty-first-century}. \textit{Locator:} Primary publisher, working-paper, or author record: bibliographic information and abstract; full-text theorem verification is not claimed. \textit{Role:} Conceptual industrial-policy framework; not a proof of optimal targeting.
\item Ann Harrison and Andrés Rodríguez-Clare (2010). \emph{Trade, Foreign Investment, and Industrial Policy for Developing Countries}. Handbook of Development Economics 5, chapter 63, 4039--4214. \url{https://www.sciencedirect.com/science/article/pii/B978044452944200001X}. \textit{Locator:} Primary publisher, working-paper, or author record: bibliographic information and abstract; full-text theorem verification is not claimed. \textit{Role:} Review of industrial-policy evidence and rationales.
\item Réka Juhász, Nathan Lane, and Dani Rodrik (2024). \emph{The New Economics of Industrial Policy}. Annual Review of Economics 16, 213--242. \url{https://doi.org/10.1146/annurev-economics-081023-024638}. \textit{Locator:} Primary publisher, working-paper, or author record: bibliographic information and abstract; full-text theorem verification is not claimed. \textit{Role:} Recent review of measurement, identification, and policy debates.
\end{enumerate}

\subsection*{Common conventions: Source mathematical conventions}

Unless an entry states otherwise, agent and firm sets are finite. State and action spaces are measurable, strategies and outcome maps are measurable, probability laws are specified, and expectations exist for the targets under discussion. Infinite-horizon discount factors lie in $(0,1)$ and discounted objectives are finite. Norms, tolerances and complexity input sizes must be specified for computational comparisons. Constants or primitives that appear locally are inputs, not empirical facts asserted by this catalogue.

\textbf{Identification.} A statistical model $m\in\mathcal M$ implies an observable law $P_m$ and target $\theta(m)$. At law $P$, its identified set is
\[
\Theta(P;\mathcal M)=\{\theta(m):m\in\mathcal M,\ P_m=P\}.
\]
Point identification holds when this set is a singleton, or equivalently when $P_{m_1}=P_{m_2}$ implies $\theta(m_1)=\theta(m_2)$ for all compatible models. Sharp bounds mean bounds attained, or approached when the boundary is not attained, by compatible models. Writing this set is a definition, not a solution: the substantive task is to establish informative restrictions and derive or estimate the set. An unrestricted-confounding model may give only trivial bounds.

\textbf{Inference.} Identification, estimability and valid uncertainty quantification are separate questions. Fix $\alpha\in(0,1)$. A confidence procedure $C_n$ over a declared law class $\mathcal P$ has asymptotic uniform coverage at level $1-\alpha$ when
\[
\liminf_{n\to\infty}\inf_{P\in\mathcal P}\Pr_P\{\theta(P)\in C_n\}\geq1-\alpha.
\]
For set-identified targets the coverage event must be defined separately, for example coverage of the full identified set. If an entry uses identification-indexed models instead of uniquely indexed laws, the same coverage convention is applied over those models. Pointwise limits do not establish uniform validity, and asymptotic validity does not guarantee finite-sample precision. Sample sizes, dependence structures and weak-identification sequences are part of each application.

\textbf{Counterfactuals.} $Y_i(a)$ is a potential outcome under a specified intervention, including its timing, implementation and versions. It is not $Y_i$ conditional on voluntarily choosing $a$. A full intervention vector permits interference; shorthand scalar treatment notation does not silently impose no interference. Assignment, selection, attrition, anticipation and measurement must be specified. A claim about an instrument's local effect is not automatically a population policy effect.

\textbf{Economic equilibrium.} A counterfactual model includes best responses, constraints, feasibility, market clearing, state transitions and expectations consistent with the stated information. When several equilibria exist, either a declared selector is an input or the target is set-valued. Comparisons across policy regimes must use the stated selection convention. Reduced-form relationships are not presumed invariant after an equilibrium-changing intervention.

\textbf{Optimization and welfare.} Optimization is over the stated feasible set. If attainment is not established, $\sup$ or $\inf$ and $\varepsilon$-optimal choices replace an asserted maximizer or minimizer. For example, with value $V=\sup_{a\in\mathcal A}W(a)<\infty$, an $\varepsilon$-optimal set is $\{a\in\mathcal A:W(a)\geq V-\varepsilon\}$ for $\varepsilon>0$. Benefits, costs, risk and health or environmental outcomes can be aggregated only using declared units and conversion weights. Interpersonal utility weights, the treatment of fiscal transfers and foreign profits, discounting, and the behavioral welfare criterion are explicit inputs. A comparative-static derivative additionally requires existence and appropriate differentiability of the selected counterfactual.

\textbf{Sources.} Local labels [S1], [S2], and so forth restart in every question. Locators identify the relevant abstract, model, section or result at the level actually checked; a bibliographic or abstract-level verification is not represented as inspection of an entire proof. Working papers are distinguished from later publications and changing versions. Source summaries are paraphrased. All URL checks and editorial formulations refer to the audit date stated above.
% END ECONOMICS STATEMENT
\end{document}
